% !TEX root = ../nc-cubic.tex
\section{楕円因子の window と比較核の基礎}
\label{cf:sec:elliptic-foundations}

分離型の直接比較に必要な楕円因子の規約を固定する。
右 twisted homogeneous coordinate ring の内部 shift、
rotation の共役式、体上の Fourier--Mukai intertwiner の存在判定と、
window1 における剰余体の像を同じ規約で記述する。
ここで扱うのは体上の因子比較であり、族の上の核 gerbe や full 残余核の降下は含めない。
楕円曲線上の twist は評価写像の cone
\[
 T_{\mathcal O_E}(F)=\operatorname{Cone}\bigl(
 R\operatorname{Hom}(\mathcal O_E,F)\otimes\mathcal O_E\longrightarrow F\bigr)
\]
で定める。

\subsection{平行移動・線束と rotation の共役}

楕円曲線の点 $b$ を次数零線束 $P_b=\mathcal O_E(b-0)$ と同一視する。
以下の式はまず代数閉体上で証明し、必要な等式を忠実平坦降下して用いる。
右加群の twisted homogeneous coordinate ring の規約は
\cite[\S3, equation~(3.1)]{NFMStaffordVdB}に合わせる。
すなわち $\sigma$ が点加群の shift を表すとき、
coherent sheaf 上の次数 shift は
\begin{equation}
\label{eq:nfm-right-thcr-shift}
 s=\sigma_*\circ(-\otimes L)
\end{equation}
である。特に $s(k(x))=k(\sigma(x))$ となる。
$\sigma^*$ と取り違えると、以下の平行移動の符号が逆になる。

\begin{lemma}[楕円曲線の rotation の共役式]
\label{lem:nfm-elliptic-rotation-conjugacy}
$E$ を標数零の代数閉体上の楕円曲線、$L$ を次数三の線束とする。
\[
 U_{a,b}=t_a^*\circ(-\otimes P_b),\qquad
 \rho=T_{\mathcal O_E}\circ(-\otimes L)
\]
と置く。関手の同型類について
\begin{equation}
\label{eq:nfm-elliptic-rotation-conjugacy}
 \rho U_{a,b}\rho^{-1}\simeq U_{-2a-b,\,3a+b}
\end{equation}
が成り立つ。
$\sigma=t_p$ とし、右加群の規約に従って
$\rho_p=T_{\mathcal O_E}\circ\sigma_*\circ(-\otimes L)$ と置くと、
\begin{equation}
\label{eq:nfm-elliptic-point-intertwiner}
 U_{a,b}\rho_pU_{a,b}^{-1}\simeq\rho
 \quad\Longleftrightarrow\quad
 3a=0,\qquad b=p.
\end{equation}
特に $\rho$ と可換する $U_{a,b}$ は $U_{a,0}$、$a\in E[3]$ に限られ、
これらは次数零 Picard 類を変えない。
\end{lemma}

\begin{proof}
平行移動は次数零 Picard 類に自明に作用するので、
$U_{a,b}$ の同型類の積は $(a,b)$ の加法である。
$t_a^*L\otimes L^{-1}\simeq P_{-3a}$ より
\[
 \operatorname{Ad}_{-\otimes L}(a,b)=(a,b+3a).
\]
また
\[
 \operatorname{Ad}_{T_{\mathcal O_E}}(a,b)=(a-b,b)
\]
である。後式の graph と線束を確認する。
twist の Fourier--Mukai 核は
$\operatorname{Cone}(\mathcal O_{E\times E}\to\mathcal O_\Delta)
=\mathcal O_{E\times E}(-\Delta)[1]$ であり、
$T_{\mathcal O_E}^{-1}k(x)=\mathcal O_E(x)$ となる。
従って共役 $T_{\mathcal O_E}U_{a,b}T_{\mathcal O_E}^{-1}$ は
$k(x)$ を $k(x-a+b)$ へ送り、$\mathcal O_E$ を $P_b$ へ送る。
点層を shift なしの点層へ送る同値の核は、その台である graph 上の線束である。
この場合の graph は $t_{a-b}$ の pullback、線束は $P_b$ なので、所要の式を得る。
これは上記の twist 核を使った畳込み計算とも一致する。

二つの共役作用の積は
\[
 A=\begin{pmatrix}-2&-1\\3&1\end{pmatrix}
\]
であり、\eqref{eq:nfm-elliptic-rotation-conjugacy}を与える。
平行移動は $\mathcal O_E$ を保つので、その pullback は
$T_{\mathcal O_E}$ と可換する。
従って $\rho_p=U_{-p,0}\rho$ であり、intertwining の条件は
\[
 (I-A)\binom ab=\binom p0,
 \qquad\text{すなわち}\qquad
 3a+b=p,\quad -3a=0
\]
に等しい。これで全ての主張が従う。
\end{proof}

\subsection{全楕円同値の分類と二段階の障害}

次に、平行移動と線束の tensor 以外も含む全ての Fourier--Mukai 同値を扱う。
$K$ を標数零の体、$E,C$ を $K$ 上の滑らかな射影 genus-one 曲線とし、
実際の次数三線束 $L_E,L_C$ を固定する。原点や $K$-有理点は仮定しない。
$p\in J_E(K)=\operatorname{Pic}^0(E)(K)$ に対し
\[
 \rho_{E,p}=T_{\mathcal O_E}\circ(t_p)_*\circ(-\otimes L_E),
 \qquad
 \rho_C=T_{\mathcal O_C}\circ(-\otimes L_C)
\]
と置く。
以下で intertwiner というときは、実際の perfect Fourier--Mukai 核による
$K$-線型同値 $F$ に、$F\rho_{E,p}\simeq\rho_CF$ という
\emph{同型の存在条件}を課す。その同型自体の標識と指定周期との整合は、
最後に別に扱う。

\begin{lemma}[数値 centralizer と graph の認識]
\label{lem:nfm-elliptic-numerical-graph}
幾何的数値格子を $(\operatorname{rank},\deg)$ によって $\mathbf Z^2$ と同定する。
$\rho_{E,p}$ と $\rho_C$ の行列はともに
\[
 R=\begin{pmatrix}-2&-1\\3&1\end{pmatrix},\qquad R^3=I
\]
であり、
\begin{equation}\label{eq:nfm-elliptic-numerical-centralizer}
 \operatorname{Cent}_{\operatorname{SL}_2(\mathbf Z)}(R)
   =\{\pm I,\pm R,\pm R^2\}.
\end{equation}
また、実際の $K$-線型 Fourier--Mukai 同値
$F:D^b(E)\to D^b(C)$ の幾何的数値作用が恒等なら、
$K$-同型 $\phi:E\xrightarrow{\sim}C$、$C$ 上の次数零線束 $M$、
整数 $a$ があって
\begin{equation}\label{eq:nfm-elliptic-numerical-graph}
 F\simeq(-\otimes M)\circ\phi_*[2a]
\end{equation}
となる。これらは同型を除いて一意であり、
$F(\mathcal O_E)$ が shift のない次数零線束なら $a=0$ である。
\end{lemma}

\begin{proof}
次数三線束の tensor と $T_{\mathcal O}$ の数値行列は、それぞれ
\[
 \begin{pmatrix}1&0\\3&1\end{pmatrix},\qquad
 \begin{pmatrix}1&-1\\0&1\end{pmatrix}
\]
であり、平行移動の数値作用は恒等である。従ってその積は $R$ となる。
$B=\left(\begin{smallmatrix}a&b\\c&d\end{smallmatrix}\right)$ とすると、
$BR=RB$ は $c=-3b$、$d=a-3b$ に等しい。
さらに
\[
 1=\det B=a^2-3ab+3b^2
   =\left(a-\frac32b\right)^2+\frac34b^2
\]
なので $|b|\le1$ である。
$b=0$ なら $a=\pm1$、$b=1$ なら $a=1,2$、
$b=-1$ なら $a=-1,-2$ であり、これが表示した六行列を与える。

graph の主張をまず代数閉包上で証明する。
滑らかな曲線では $\operatorname{Ext}^{\ge2}=0$ なので、
有界複体は cohomology 層の shift の直和である。
$F(k(x))$ は indecomposable で数値類 $(0,1)$ を持つため、
ある点 $\phi(x)$ と整数 $a(x)$ に対し
$F(k(x))\simeq k(\phi(x))[2a(x)]$ となる。
kernel の第一射影に関する relative perfect 性と Tor amplitude の開性から
$a(x)$ は局所定数であり、$E$ の連結性によって一定である。
shift を除いた kernel は第一射影について平坦な長さ一の層の族なので、
一つの graph 上の線束となる。
逆 kernel にも同じ議論を適用すると、graph の射は同型である。
これで\eqref{eq:nfm-elliptic-numerical-graph}を得る。
数値作用が恒等なので $\deg M=0$ であり、
$F(\mathcal O_E)=M[2a]$ から最後の主張も従う。
この graph 認識は
\cite[Proposition~4.2, Corollary~4.3 and Theorem~5.3]{NFMHilleVdBFM}
の証明と同じ議論である。

$K$ 上への降下では、kernel が実際の $K$ 上の対象であることを用いる。
忠実平坦な体拡大後に cohomology が一つの次数に集中するので、
$K$ 上でも同じ次数の層となる。
その平坦性、graph 上の可逆性は fpqc 局所的性質である。
scheme 理論的 support の $E$ への射は代数閉包上で同型なので、
$K$ 上でも同型となり、$\phi$ と線束が降下する。
graph とその上の線束の一意性から、表示の一意性も従う。
ここで $K$ 上の数値格子が次数一の対象を含むとは仮定していない。
\end{proof}

\begin{theorem}[構造層の像で規格化した全楕円 intertwiner]
\label{thm:nfm-elliptic-normalized-intertwiners}
$F:D^b(E)\to D^b(C)$ を実際の $K$-線型 Fourier--Mukai 同値とし、
$F(\mathcal O_E)$ は shift のない次数零線束とする。
このとき $F\rho_{E,p}\simeq\rho_CF$ であることは、
\begin{equation}\label{eq:nfm-elliptic-normalized-intertwiners}
 F\simeq(-\otimes M)\circ\phi_*,\qquad
 \phi^*L_C\simeq L_E,\qquad
 [M]=\operatorname{Jac}(\phi)(p)
\end{equation}
を満たす $K$-同型 $\phi:E\to C$ と実線束 $M$ があることと同値である。
ここで $\operatorname{Jac}(\phi)$ は $\phi$ による共変な Picard 同型である。
\end{theorem}

\begin{proof}
$F$ の幾何的数値行列 $B$ は $e_1=(1,0)^{\mathsf t}$ を固定し、
$BR=RB$ を満たす。$e_1$ と $Re_1=(-2,3)^{\mathsf t}$ は
$\mathbf Q$ 上独立なので $B=I$ である。
前補題によって $F=(-\otimes M)\phi_*$ と書ける。
\[
 q=\operatorname{Jac}(\phi)(p),\qquad
 L'=\phi_*L_E,\qquad d=[L'\otimes L_C^{-1}],\qquad m=[M]
\]
と置く。
代数閉包上で原点を取り、補題\ref{lem:nfm-elliptic-rotation-conjugacy}の
$U_{a,b}$ を用いると
\[
 \phi_*\rho_{E,p}\phi_*^{-1}=U_{-q,0}\rho_CU_{0,d}
\]
である。従って同補題の共役式から
\begin{equation}\label{eq:nfm-elliptic-graph-conjugation}
 \begin{split}
 F\rho_{E,p}F^{-1}
 &=U_{0,m}U_{-q,0}\rho_CU_{0,d}U_{0,-m}\\
 &=U_{m-q-d,d}\rho_C
 \end{split}
\end{equation}
となる。
$U_{a,b}\simeq\mathrm{id}$ なら、$\mathcal O$ の像から $b=0$、
点層の像から $a=0$ である。
従って intertwining は $d=0$、$m=q$ に等しい。
特に非偏極的な $\phi$ のずれを線束 $M$ で補償することはできない。

これらは相対 Picard 類の等式なので $K$ 上へ降りる。
$d=0$ は、実線束 $L'\otimes L_C^{-1}$ が
$\operatorname{Spec}K$ からの線束であることを意味し、従って自明である。
逆向きの functor の同型も $K$ 上で存在する。
実際、比較する二つの同値 kernel は代数閉包上で同型であり、
その $\operatorname{Hom}^0$ は $\mathrm{HH}^0(E_{\overline K})=\overline K$
によって一次元である。
proper perfect Hom の体拡大により $K$ 上でも一次元であり、
その非零元は代数閉包上で同型となるため、既に $K$ 上で同型である。
\end{proof}

$p\in J_E(K)$ に対し、$p$ を表す実際の次数零線束の Picard gerbe の類を
$\delta_E(p)\in\operatorname{Br}(K)$ と書く。
その消滅は、$p$ を表す線束が $E$ 上に存在することを意味する。

\begin{corollary}[全数値型を含む active 同値の存在判定]
\label{cor:nfm-elliptic-any-intertwiner-criterion}
$F(\mathcal O_E)$ に条件を課さない場合にも、
実際の $K$-線型 Fourier--Mukai intertwiner
$(D^b(E),\rho_{E,p})\to(D^b(C),\rho_C)$ が存在することは、
次の二条件と同値である：
\begin{enumerate}
 \item $(E,L_E)$ と $(C,L_C)$ の $K$ 上の偏極曲線同型が存在する。
 \item $\delta_E(p)=0$、すなわち $p$ は $E$ 上の実際の次数零線束で表される。
\end{enumerate}
特に同じ $(E,L)$ 上では
\begin{equation}\label{eq:nfm-elliptic-any-intertwiner-obstruction}
 \{F:F\rho_{E,p}\simeq\rho_EF\}\ne\varnothing
 \quad\Longleftrightarrow\quad \delta_E(p)=0.
\end{equation}
この主張は intertwining 同型自体の周期付き rigidification を含まない。
\end{corollary}

\begin{proof}
任意の intertwiner の数値行列は
\eqref{eq:nfm-elliptic-numerical-centralizer}に属する。
$K$ 上で定義された $\rho_C$ の冪と shift を postcompose して、
数値作用を恒等にする。
補題\ref{lem:nfm-elliptic-numerical-graph}によって graph と
次数零線束の偶 shift を得るので、その shift も除けば前定理を適用できる。
従って偏極曲線同型 $\phi$ と
$\operatorname{Jac}(\phi)(p)$ を表す実線束を得る。
実際の $K$-同型 $\phi$ による線束の移送は Picard gerbe の band を保つから
\[
 \delta_C(\operatorname{Jac}(\phi)(p))=\delta_E(p)
\]
であり、必要性を得る。
逆に $p$ の実線束を $\phi$ で移送すれば、前定理から所要の同値を得る。
\end{proof}

\subsection{右座標環の window と内部 shift}

以下の window の議論では基礎体を $k$ と書き、
$\kappa_A$ を増大剰余体 $k=A/A_{\ge1}$ の特異圏における像とする。

\begin{lemma}[右楕円座標環の window1 と内部 shift]
\label{lem:nfm-twisted-elliptic-window}
$A=B(E,L,\sigma)$ を右 twisted homogeneous coordinate ring とする。
$\operatorname{qgr}A\simeq\operatorname{coh}E$ の下で
$\pi A\leftrightarrow\mathcal O_E$、$(1)\leftrightarrow s$ は
\eqref{eq:nfm-right-thcr-shift}に従うものとする。
$A$ に Gorenstein parameter 零の Orlov 同値が適用でき、
window1 を $\Phi_1=q_{\mathrm{sg}}R\Gamma_{\ge1}$ と取る。
このとき
\begin{equation}
\label{eq:nfm-twisted-active-window}
 \kappa_A\longleftrightarrow\mathcal O_E[1],\qquad
 \Phi_1^{-1}(1)\Phi_1=T_{\mathcal O_E}\circ s.
\end{equation}
\end{lemma}

\begin{proof}
この window の存在については
\cite[Theorem~2.5(iii)]{NFMOrlovSingularities}を用いる。
$m>0$ に対して $s^m\mathcal O_E$ は次数 $3m$ の線束なので、
その高次 cohomology は消える。
従って $R\Gamma_{\ge1}\mathcal O_E=A_{\ge1}$ であり、
$0\to A_{\ge1}\to A\to k\to0$ を特異圏へ写すと第一式を得る。
任意の $F$ に対し
\[
 (R\Gamma_{\ge1}F)(1)=R\Gamma_{\ge0}(sF)
\]
であり、次数零の切断は三角形
\[
 R\Gamma_{\ge1}(sF)\longrightarrow R\Gamma_{\ge0}(sF)
 \longrightarrow R\operatorname{Hom}(\mathcal O_E,sF)\otimes k
 \longrightarrow
\]
を与える。
特異圏へ写し、第一式を用いて三角形を回転すると、評価写像の三角形
\[
 R\operatorname{Hom}(\mathcal O_E,sF)\otimes\mathcal O_E
 \longrightarrow sF\longrightarrow\Phi_1^{-1}(1)\Phi_1(F)
 \longrightarrow
\]
を得る。この cone は $T_{\mathcal O_E}(sF)$ である。
従って剰余体と内部 shift の双方について同じ window の規約を得た。
この切断の議論は $A$ の可換性を用いない。
\end{proof}

